% Pratinjau gambar TikZ skripsi v3.0 -- satu gambar per halaman.
% Kelas dan preambul sama dengan praskripsi.tex. Opsi `ketat` DIHAPUS di sini
% supaya pratinjau tetap bisa dikompilasi tanpa font Times New Roman/Arial.
\documentclass[jenis=praskripsi]{upnjatim-skripsi}
\upnsetup{
  % CEK: Judul sementara, belum final. Diputuskan bersama dosen pembimbing.
  % Istilah asing (time frame, hyperparameter) belum dicetak miring di sampul
  % karena nilai kunci ini juga dipakai sebagai pdftitle.
  judul={Prediksi Harga BTC/IDR pada Berbagai Time Frame dengan Model Boosting dan Optimasi Hyperparameter},
  nama={Ferdi Endahas Ahmad},
  npm={23081010165},
  pembimbing-satu={Dr. Rizky Parlika, S.Kom., M.Kom.},
  pembimbing-dua={Angga Lisdiyanto, S.Kom., M.Kom.},
  program-studi={Informatika},
  tahun={2026}
}


% Library TikZ tambahan yang belum dimuat di upnjatim-skripsi.cls.
% Keduanya perlu dipindahkan ke .cls sebelum gambar dimasukkan ke chapters/.
\usetikzlibrary{patterns}                  % arsir untuk purging dan jeda
\usetikzlibrary{decorations.pathreplacing} % kurung kurawal jendela fitur

\begin{document}
\pagestyle{empty}

\newcommand{\pratinjau}[2]{%
  \begin{center}
    \textbf{#1}\par\medskip
    \input{#2}
  \end{center}
  \clearpage}

\pratinjau{G2-1 \quad \texttt{g2-horizon-naive}}{g2-horizon-naive}
\pratinjau{G2-2 \quad \texttt{g2-gradient-boosting}}{g2-gradient-boosting}
\pratinjau{G2-3 \quad \texttt{g2-leafwise-levelwise}}{g2-leafwise-levelwise}
\pratinjau{G2-4 \quad \texttt{g2-pohon-simetris}}{g2-pohon-simetris}
\pratinjau{G2-5 \quad \texttt{g2-adaboost-r2}}{g2-adaboost-r2}
\pratinjau{G2-6 \quad \texttt{g2-grid-random}}{g2-grid-random}
\pratinjau{G2-7 \quad \texttt{g2-kerangka-pemikiran}}{g2-kerangka-pemikiran}
\pratinjau{G3-1 \quad \texttt{g3-alur-penelitian}}{g3-alur-penelitian}
\pratinjau{G3-2 \quad \texttt{g3-pengumpulan-data}}{g3-pengumpulan-data}
\pratinjau{G3-3 \quad \texttt{g3-fitur-target}}{g3-fitur-target}
\pratinjau{G3-4 \quad \texttt{g3-split-purging}}{g3-split-purging}
\pratinjau{G3-5 \quad \texttt{g3-walk-forward}}{g3-walk-forward}
\pratinjau{G3-6 \quad \texttt{g3-tahapan-eksperimen}}{g3-tahapan-eksperimen}

\end{document}
